
\documentclass{article}
\usepackage{colortbl}
\usepackage{makecell}
\usepackage{multirow}
\usepackage{supertabular}

\begin{document}

\newcounter{utterance}

\centering \large Interaction Transcript for game `matchit\_ascii', experiment `different\_grid', episode 1 with glm{-}5.2.
\vspace{24pt}

{ \footnotesize  \setcounter{utterance}{1}
\setlength{\tabcolsep}{0pt}
\begin{supertabular}{c@{$\;$}|p{.15\linewidth}@{}p{.15\linewidth}p{.15\linewidth}p{.15\linewidth}p{.15\linewidth}p{.15\linewidth}}
    \# & $\;$A & \multicolumn{4}{c}{Game Master} & $\;\:$B\\
    \hline

    \theutterance \stepcounter{utterance}  
    & & \multicolumn{4}{p{0.6\linewidth}}{
        \cellcolor[rgb]{0.9,0.9,0.9}{
            \makecell[{{p{\linewidth}}}]{
                \texttt{\tiny{[P1$\langle$GM]}}
                \texttt{Tá tú ag glacadh páirte i gcluiche tomhais comhoibríoch.} \\
\texttt{Is é an sprioc a fháil amach an bhfuil an greille seo agus pictiúr eile nach féidir liomsa amháin a fheiceáil, mar an gcéanna.} \\
\texttt{Déan cur síos ar do ghreille ar dtús, le do thoil. Ansin, tabharfaidh mé mo chur síos agus is féidir linn ceisteanna a chur ar a chéile faoi na greillí chun a dhéanamh amach an bhfuil siad mar an gcéanna. Anois cuir tús le do chur síos gairid ar an ngreille le "TUAIRISC:" agus an cur síos ina dhiaidh. Ná cuir aon rud eile leis. Seo é do ghreille:} \\
\texttt{X X X X X} \\
\texttt{▢ ▢ ▢ ▢ X} \\
\texttt{▢ ▢ X X X} \\
\texttt{▢ ▢ ▢ ▢ X} \\
\texttt{X X X X X} \\
            }
        }
    }
    & \\ \\

    \theutterance \stepcounter{utterance}  
    & \multicolumn{4}{p{0.6\linewidth}}{
        \cellcolor[rgb]{0.9,0.9,0.9}{
            \makecell[{{p{\linewidth}}}]{
                \texttt{\tiny{[P1$\rangle$GM]}}
                \texttt{TUAIRISC: Greille 5x5. Is iad na cearnóga i líne uachtarach ná: X, X, X, X, X. Is iad na cearnóga i líne 2 ná: ▢, ▢, ▢, ▢, X. Is iad na cearnóga i líne 3 ná: ▢, ▢, X, X, X. Is iad na cearnóga i líne 4 ná: ▢, ▢, ▢, ▢, X. Is iad na cearnóga i líne íochtarach ná: X, X, X, X, X.} \\
            }
        }
    }
    & & \\ \\

    \theutterance \stepcounter{utterance}  
    & & & \multicolumn{2}{p{0.3\linewidth}}{
        \cellcolor[rgb]{0.9,0.9,0.9}{
            \makecell[{{p{\linewidth}}}]{
                \texttt{\tiny{[GM$|$GM]}}
                \texttt{continue} \\
            }
        }
    }
    & & \\ \\

    \theutterance \stepcounter{utterance}  
    & & \multicolumn{4}{p{0.6\linewidth}}{
        \cellcolor[rgb]{0.9,0.9,0.9}{
            \makecell[{{p{\linewidth}}}]{
                \texttt{\tiny{[GM$\rangle$P2]}}
                \texttt{Tá tú ag glacadh páirte i gcluiche tomhais comhoibríoch.} \\
\texttt{Is é an sprioc a fháil amach an bhfuil an greille seo agus pictiúr eile nach féidir liomsa amháin a fheiceáil, mar an gcéanna.} \\
\texttt{Déan cur síos ar do ghreille ar dtús, le do thoil. Ansin, tabharfaidh mé mo chur síos agus is féidir linn ceisteanna a chur ar a chéile faoi na greillí chun a dhéanamh amach an bhfuil siad mar an gcéanna. Anois cuir tús le do chur síos gairid ar an ngreille le "TUAIRISC:" agus an cur síos ina dhiaidh. Ná cuir aon rud eile leis. Seo é do ghreille:} \\
\texttt{X ▢ ▢ ▢ X} \\
\texttt{X ▢ ▢ X X} \\
\texttt{X ▢ X ▢ X} \\
\texttt{X X ▢ ▢ X} \\
\texttt{X ▢ ▢ ▢ X} \\
            }
        }
    }
    & \\ \\

    \theutterance \stepcounter{utterance}  
    & & & \multicolumn{4}{p{0.6\linewidth}}{
        \cellcolor[rgb]{0.9,0.9,0.9}{
            \makecell[{{p{\linewidth}}}]{
                \texttt{\tiny{[GM$\langle$P2]}}
                \texttt{TUAIRISC: Cúig fhadlíne atá sa ghreille, cúig cheallóg i ngach fhadlíne. Tá 'X' sa chéad agus sa cheann deireanigh ceallóg i ngach fhadlíne amháin, ag cruthú fráma nó imlíne sheachtrach. Sa chéad fhadlíne, níl aon 'X' sa lar. Sa dara fhadlíne, tá 'X' sa cheathrú ceallóg. Sa tríú fhadlíne, tá 'X' sa tríú ceallóg. Sa cheathrú fhadlíne, tá 'X' sa dara agus sa tríú ceallóg. Sa chúigiú fhadlíne, níl aon 'X' sa lár.} \\
            }
        }
    }
     \\ \\

    \theutterance \stepcounter{utterance}  
    & & & \multicolumn{2}{p{0.3\linewidth}}{
        \cellcolor[rgb]{0.9,0.9,0.9}{
            \makecell[{{p{\linewidth}}}]{
                \texttt{\tiny{[GM$|$GM]}}
                \texttt{continue} \\
            }
        }
    }
    & & \\ \\

    \theutterance \stepcounter{utterance}  
    & & \multicolumn{4}{p{0.6\linewidth}}{
        \cellcolor[rgb]{0.9,0.9,0.9}{
            \makecell[{{p{\linewidth}}}]{
                \texttt{\tiny{[GM$\rangle$P2]}}
                \texttt{Seo é mo CHUR SÍOSTUAIRISC: Greille 5x5. Is iad na cearnóga i líne uachtarach ná: X, X, X, X, X. Is iad na cearnóga i líne 2 ná: ▢, ▢, ▢, ▢, X. Is iad na cearnóga i líne 3 ná: ▢, ▢, X, X, X. Is iad na cearnóga i líne 4 ná: ▢, ▢, ▢, ▢, X. Is iad na cearnóga i líne íochtarach ná: X, X, X, X, X.} \\
\texttt{Anois cuir ceist chun gnéithe nua de mo ghreille a fháil amach a d'fhéadfadh a bheith difriúil ó do ghreille. Tosaigh le "CEIST:" agus ná cuir aon rud eile leis.} \\
            }
        }
    }
    & \\ \\

    \theutterance \stepcounter{utterance}  
    & & & \multicolumn{4}{p{0.6\linewidth}}{
        \cellcolor[rgb]{0.9,0.9,0.9}{
            \makecell[{{p{\linewidth}}}]{
                \texttt{\tiny{[GM$\langle$P2]}}
                \texttt{CEIST: An bhfuil 'X' ar aiséime nó i lár do ghreille?} \\
            }
        }
    }
     \\ \\

    \theutterance \stepcounter{utterance}  
    & & & \multicolumn{2}{p{0.3\linewidth}}{
        \cellcolor[rgb]{0.9,0.9,0.9}{
            \makecell[{{p{\linewidth}}}]{
                \texttt{\tiny{[GM$|$GM]}}
                \texttt{continue} \\
            }
        }
    }
    & & \\ \\

    \theutterance \stepcounter{utterance}  
    & & \multicolumn{4}{p{0.6\linewidth}}{
        \cellcolor[rgb]{0.9,0.9,0.9}{
            \makecell[{{p{\linewidth}}}]{
                \texttt{\tiny{[P1$\langle$GM]}}
                \texttt{Seo é mo CHUR SÍOSTUAIRISC: Cúig fhadlíne atá sa ghreille, cúig cheallóg i ngach fhadlíne. Tá 'X' sa chéad agus sa cheann deireanigh ceallóg i ngach fhadlíne amháin, ag cruthú fráma nó imlíne sheachtrach. Sa chéad fhadlíne, níl aon 'X' sa lar. Sa dara fhadlíne, tá 'X' sa cheathrú ceallóg. Sa tríú fhadlíne, tá 'X' sa tríú ceallóg. Sa cheathrú fhadlíne, tá 'X' sa dara agus sa tríú ceallóg. Sa chúigiú fhadlíne, níl aon 'X' sa lár.} \\
\texttt{CEIST: An bhfuil 'X' ar aiséime nó i lár do ghreille?Tosaigh do fhreagra le "FREAGRA:" agus ná cuir aon rud eile leis.} \\
            }
        }
    }
    & \\ \\

    \theutterance \stepcounter{utterance}  
    & \multicolumn{4}{p{0.6\linewidth}}{
        \cellcolor[rgb]{0.9,0.9,0.9}{
            \makecell[{{p{\linewidth}}}]{
                \texttt{\tiny{[P1$\rangle$GM]}}
                \texttt{FREAGRA: Níl. Níl aon 'X' sa cheallóg lár (lár na greille) agus tá an ghreille aiséimeach.} \\
            }
        }
    }
    & & \\ \\

    \theutterance \stepcounter{utterance}  
    & & & \multicolumn{2}{p{0.3\linewidth}}{
        \cellcolor[rgb]{0.9,0.9,0.9}{
            \makecell[{{p{\linewidth}}}]{
                \texttt{\tiny{[GM$|$GM]}}
                \texttt{continue} \\
            }
        }
    }
    & & \\ \\

    \theutterance \stepcounter{utterance}  
    & & \multicolumn{4}{p{0.6\linewidth}}{
        \cellcolor[rgb]{0.9,0.9,0.9}{
            \makecell[{{p{\linewidth}}}]{
                \texttt{\tiny{[P1$\langle$GM]}}
                \texttt{Anois cuir ceist chun gnéithe nua de mo ghreille a fháil amach a d'fhéadfadh a bheith difriúil ó do ghreille. Tosaigh le "CEIST:" agus ná cuir aon rud eile leis.} \\
            }
        }
    }
    & \\ \\

    \theutterance \stepcounter{utterance}  
    & \multicolumn{4}{p{0.6\linewidth}}{
        \cellcolor[rgb]{0.9,0.9,0.9}{
            \makecell[{{p{\linewidth}}}]{
                \texttt{\tiny{[P1$\rangle$GM]}}
                \texttt{CEIST: Sa cheathrú fhadlíne, an bhfuil 'X' sa dara agus sa tríú ceallóg, agus an bhfuil na cealla eile sa líne sin folamh (▢)?} \\
            }
        }
    }
    & & \\ \\

    \theutterance \stepcounter{utterance}  
    & & & \multicolumn{2}{p{0.3\linewidth}}{
        \cellcolor[rgb]{0.9,0.9,0.9}{
            \makecell[{{p{\linewidth}}}]{
                \texttt{\tiny{[GM$|$GM]}}
                \texttt{continue} \\
            }
        }
    }
    & & \\ \\

    \theutterance \stepcounter{utterance}  
    & & \multicolumn{4}{p{0.6\linewidth}}{
        \cellcolor[rgb]{0.9,0.9,0.9}{
            \makecell[{{p{\linewidth}}}]{
                \texttt{\tiny{[GM$\rangle$P2]}}
                \texttt{FREAGRA: Níl. Níl aon 'X' sa cheallóg lár (lár na greille) agus tá an ghreille aiséimeach.} \\
\texttt{CEIST: Sa cheathrú fhadlíne, an bhfuil 'X' sa dara agus sa tríú ceallóg, agus an bhfuil na cealla eile sa líne sin folamh (▢)?Tosaigh do fhreagra le "FREAGRA:" agus ná cuir aon rud eile leis.} \\
            }
        }
    }
    & \\ \\

    \theutterance \stepcounter{utterance}  
    & & & \multicolumn{4}{p{0.6\linewidth}}{
        \cellcolor[rgb]{0.9,0.9,0.9}{
            \makecell[{{p{\linewidth}}}]{
                \texttt{\tiny{[GM$\langle$P2]}}
                \texttt{FREAGRA: Ní hea. Sa cheathrú fhadlíne níl 'X' sa dara ná sa tríú ceallóg, tá siad folamh (▢). Níl ann ach ceallóg amháin 'X' sa cheathrú fhadlíne, agus sin í sa chúigiú ceallóg.} \\
            }
        }
    }
     \\ \\

    \theutterance \stepcounter{utterance}  
    & & & \multicolumn{2}{p{0.3\linewidth}}{
        \cellcolor[rgb]{0.9,0.9,0.9}{
            \makecell[{{p{\linewidth}}}]{
                \texttt{\tiny{[GM$|$GM]}}
                \texttt{continue} \\
            }
        }
    }
    & & \\ \\

    \theutterance \stepcounter{utterance}  
    & & \multicolumn{4}{p{0.6\linewidth}}{
        \cellcolor[rgb]{0.9,0.9,0.9}{
            \makecell[{{p{\linewidth}}}]{
                \texttt{\tiny{[GM$\rangle$P2]}}
                \texttt{Anois cuir ceist chun gnéithe nua de mo ghreille a fháil amach a d'fhéadfadh a bheith difriúil ó do ghreille. Tosaigh le "CEIST:" agus ná cuir aon rud eile leis.} \\
            }
        }
    }
    & \\ \\

    \theutterance \stepcounter{utterance}  
    & & & \multicolumn{4}{p{0.6\linewidth}}{
        \cellcolor[rgb]{0.9,0.9,0.9}{
            \makecell[{{p{\linewidth}}}]{
                \texttt{\tiny{[GM$\langle$P2]}}
                \texttt{CEIST: An bhfuil cealla 'X' agat sa chéad colún, ón bharr go dtí an bun?} \\
            }
        }
    }
     \\ \\

    \theutterance \stepcounter{utterance}  
    & & & \multicolumn{2}{p{0.3\linewidth}}{
        \cellcolor[rgb]{0.9,0.9,0.9}{
            \makecell[{{p{\linewidth}}}]{
                \texttt{\tiny{[GM$|$GM]}}
                \texttt{continue} \\
            }
        }
    }
    & & \\ \\

    \theutterance \stepcounter{utterance}  
    & & \multicolumn{4}{p{0.6\linewidth}}{
        \cellcolor[rgb]{0.9,0.9,0.9}{
            \makecell[{{p{\linewidth}}}]{
                \texttt{\tiny{[P1$\langle$GM]}}
                \texttt{FREAGRA: Ní hea. Sa cheathrú fhadlíne níl 'X' sa dara ná sa tríú ceallóg, tá siad folamh (▢). Níl ann ach ceallóg amháin 'X' sa cheathrú fhadlíne, agus sin í sa chúigiú ceallóg.} \\
\texttt{CEIST: An bhfuil cealla 'X' agat sa chéad colún, ón bharr go dtí an bun?Tosaigh do fhreagra le "FREAGRA:" agus ná cuir aon rud eile leis.} \\
            }
        }
    }
    & \\ \\

    \theutterance \stepcounter{utterance}  
    & \multicolumn{4}{p{0.6\linewidth}}{
        \cellcolor[rgb]{0.9,0.9,0.9}{
            \makecell[{{p{\linewidth}}}]{
                \texttt{\tiny{[P1$\rangle$GM]}}
                \texttt{FREAGRA: Níl. Níl 'X' sa chéad colún i líne a dó ná i líne a ceathair; tá 'X' ag an mbarr agus ag an mbun sa chéad colún amháin.} \\
            }
        }
    }
    & & \\ \\

    \theutterance \stepcounter{utterance}  
    & & & \multicolumn{2}{p{0.3\linewidth}}{
        \cellcolor[rgb]{0.9,0.9,0.9}{
            \makecell[{{p{\linewidth}}}]{
                \texttt{\tiny{[GM$|$GM]}}
                \texttt{continue} \\
            }
        }
    }
    & & \\ \\

    \theutterance \stepcounter{utterance}  
    & & \multicolumn{4}{p{0.6\linewidth}}{
        \cellcolor[rgb]{0.9,0.9,0.9}{
            \makecell[{{p{\linewidth}}}]{
                \texttt{\tiny{[P1$\langle$GM]}}
                \texttt{Anois cuir ceist chun gnéithe nua de mo ghreille a fháil amach a d'fhéadfadh a bheith difriúil ó do ghreille. Tosaigh le "CEIST:" agus ná cuir aon rud eile leis.} \\
            }
        }
    }
    & \\ \\

    \theutterance \stepcounter{utterance}  
    & \multicolumn{4}{p{0.6\linewidth}}{
        \cellcolor[rgb]{0.9,0.9,0.9}{
            \makecell[{{p{\linewidth}}}]{
                \texttt{\tiny{[P1$\rangle$GM]}}
                \texttt{CEIST: Sa dara fhadlíne, an bhfuil 'X' sa chúigiú ceallóg agus na ceithre cheallóg eile folamh (▢)?} \\
            }
        }
    }
    & & \\ \\

    \theutterance \stepcounter{utterance}  
    & & & \multicolumn{2}{p{0.3\linewidth}}{
        \cellcolor[rgb]{0.9,0.9,0.9}{
            \makecell[{{p{\linewidth}}}]{
                \texttt{\tiny{[GM$|$GM]}}
                \texttt{continue} \\
            }
        }
    }
    & & \\ \\

    \theutterance \stepcounter{utterance}  
    & & \multicolumn{4}{p{0.6\linewidth}}{
        \cellcolor[rgb]{0.9,0.9,0.9}{
            \makecell[{{p{\linewidth}}}]{
                \texttt{\tiny{[GM$\rangle$P2]}}
                \texttt{FREAGRA: Níl. Níl 'X' sa chéad colún i líne a dó ná i líne a ceathair; tá 'X' ag an mbarr agus ag an mbun sa chéad colún amháin.} \\
\texttt{CEIST: Sa dara fhadlíne, an bhfuil 'X' sa chúigiú ceallóg agus na ceithre cheallóg eile folamh (▢)?Tosaigh do fhreagra le "FREAGRA:" agus ná cuir aon rud eile leis.} \\
            }
        }
    }
    & \\ \\

    \theutterance \stepcounter{utterance}  
    & & & \multicolumn{4}{p{0.6\linewidth}}{
        \cellcolor[rgb]{0.9,0.9,0.9}{
            \makecell[{{p{\linewidth}}}]{
                \texttt{\tiny{[GM$\langle$P2]}}
                \texttt{FREAGRA: Ní hea. Sa dara fhadlíne, tá 'X' sa cheathrú agus sa chúigiú ceallóg, agus tá na chéad trí cheallóg folamh (▢).} \\
            }
        }
    }
     \\ \\

    \theutterance \stepcounter{utterance}  
    & & & \multicolumn{2}{p{0.3\linewidth}}{
        \cellcolor[rgb]{0.9,0.9,0.9}{
            \makecell[{{p{\linewidth}}}]{
                \texttt{\tiny{[GM$|$GM]}}
                \texttt{continue} \\
            }
        }
    }
    & & \\ \\

    \theutterance \stepcounter{utterance}  
    & & \multicolumn{4}{p{0.6\linewidth}}{
        \cellcolor[rgb]{0.9,0.9,0.9}{
            \makecell[{{p{\linewidth}}}]{
                \texttt{\tiny{[GM$\rangle$P2]}}
                \texttt{Anois cuir ceist chun gnéithe nua de mo ghreille a fháil amach a d'fhéadfadh a bheith difriúil ó do ghreille. Tosaigh le "CEIST:" agus ná cuir aon rud eile leis.} \\
            }
        }
    }
    & \\ \\

    \theutterance \stepcounter{utterance}  
    & & & \multicolumn{4}{p{0.6\linewidth}}{
        \cellcolor[rgb]{0.9,0.9,0.9}{
            \makecell[{{p{\linewidth}}}]{
                \texttt{\tiny{[GM$\langle$P2]}}
                \texttt{CEIST: An bhfuil 'X' sa chúigiú colún i ngach fhadlíne den ghreille?} \\
            }
        }
    }
     \\ \\

    \theutterance \stepcounter{utterance}  
    & & & \multicolumn{2}{p{0.3\linewidth}}{
        \cellcolor[rgb]{0.9,0.9,0.9}{
            \makecell[{{p{\linewidth}}}]{
                \texttt{\tiny{[GM$|$GM]}}
                \texttt{continue} \\
            }
        }
    }
    & & \\ \\

    \theutterance \stepcounter{utterance}  
    & & \multicolumn{4}{p{0.6\linewidth}}{
        \cellcolor[rgb]{0.9,0.9,0.9}{
            \makecell[{{p{\linewidth}}}]{
                \texttt{\tiny{[P1$\langle$GM]}}
                \texttt{FREAGRA: Ní hea. Sa dara fhadlíne, tá 'X' sa cheathrú agus sa chúigiú ceallóg, agus tá na chéad trí cheallóg folamh (▢).} \\
\texttt{CEIST: An bhfuil 'X' sa chúigiú colún i ngach fhadlíne den ghreille?Tosaigh do fhreagra le "FREAGRA:" agus ná cuir aon rud eile leis.} \\
            }
        }
    }
    & \\ \\

    \theutterance \stepcounter{utterance}  
    & \multicolumn{4}{p{0.6\linewidth}}{
        \cellcolor[rgb]{0.9,0.9,0.9}{
            \makecell[{{p{\linewidth}}}]{
                \texttt{\tiny{[P1$\rangle$GM]}}
                \texttt{FREAGRA: Tá. Tá 'X' sa chúigiú colún i ngach fhadlíne den ghreille.} \\
            }
        }
    }
    & & \\ \\

    \theutterance \stepcounter{utterance}  
    & & & \multicolumn{2}{p{0.3\linewidth}}{
        \cellcolor[rgb]{0.9,0.9,0.9}{
            \makecell[{{p{\linewidth}}}]{
                \texttt{\tiny{[GM$|$GM]}}
                \texttt{continue} \\
            }
        }
    }
    & & \\ \\

    \theutterance \stepcounter{utterance}  
    & & \multicolumn{4}{p{0.6\linewidth}}{
        \cellcolor[rgb]{0.9,0.9,0.9}{
            \makecell[{{p{\linewidth}}}]{
                \texttt{\tiny{[P1$\langle$GM]}}
                \texttt{Anois cuir ceist chun gnéithe nua de mo ghreille a fháil amach a d'fhéadfadh a bheith difriúil ó do ghreille. Tosaigh le "CEIST:" agus ná cuir aon rud eile leis.} \\
            }
        }
    }
    & \\ \\

    \theutterance \stepcounter{utterance}  
    & \multicolumn{4}{p{0.6\linewidth}}{
        \cellcolor[rgb]{0.9,0.9,0.9}{
            \makecell[{{p{\linewidth}}}]{
                \texttt{\tiny{[P1$\rangle$GM]}}
                \texttt{CEIST: Sa tríú fhadlíne, an bhfuil 'X' sa tríú, sa cheathrú agus sa chúigiú ceallóg, agus na chéad dhá cheallóg folamh (▢)?} \\
            }
        }
    }
    & & \\ \\

    \theutterance \stepcounter{utterance}  
    & & & \multicolumn{2}{p{0.3\linewidth}}{
        \cellcolor[rgb]{0.9,0.9,0.9}{
            \makecell[{{p{\linewidth}}}]{
                \texttt{\tiny{[GM$|$GM]}}
                \texttt{continue} \\
            }
        }
    }
    & & \\ \\

    \theutterance \stepcounter{utterance}  
    & & \multicolumn{4}{p{0.6\linewidth}}{
        \cellcolor[rgb]{0.9,0.9,0.9}{
            \makecell[{{p{\linewidth}}}]{
                \texttt{\tiny{[GM$\rangle$P2]}}
                \texttt{FREAGRA: Tá. Tá 'X' sa chúigiú colún i ngach fhadlíne den ghreille.} \\
\texttt{CEIST: Sa tríú fhadlíne, an bhfuil 'X' sa tríú, sa cheathrú agus sa chúigiú ceallóg, agus na chéad dhá cheallóg folamh (▢)?Tosaigh do fhreagra le "FREAGRA:" agus ná cuir aon rud eile leis.} \\
            }
        }
    }
    & \\ \\

    \theutterance \stepcounter{utterance}  
    & & & \multicolumn{4}{p{0.6\linewidth}}{
        \cellcolor[rgb]{0.9,0.9,0.9}{
            \makecell[{{p{\linewidth}}}]{
                \texttt{\tiny{[GM$\langle$P2]}}
                \texttt{FREAGRA: Ní hea. Sa tríú fhadlíne, tá 'X' sa chéad, sa tríú agus sa chúigiú ceallóg amháin. Tá an dara agus an ceathrú ceallóg folamh (▢).} \\
            }
        }
    }
     \\ \\

    \theutterance \stepcounter{utterance}  
    & & & \multicolumn{2}{p{0.3\linewidth}}{
        \cellcolor[rgb]{0.9,0.9,0.9}{
            \makecell[{{p{\linewidth}}}]{
                \texttt{\tiny{[GM$|$GM]}}
                \texttt{continue} \\
            }
        }
    }
    & & \\ \\

    \theutterance \stepcounter{utterance}  
    & & \multicolumn{4}{p{0.6\linewidth}}{
        \cellcolor[rgb]{0.9,0.9,0.9}{
            \makecell[{{p{\linewidth}}}]{
                \texttt{\tiny{[GM$\rangle$P2]}}
                \texttt{Anois déan cinneadh. Cad a cheapann tú: an bhfuil do ghreille agus an ghreille eile a thuairiscítear mar an gcéanna? Scríobh "CINNEADH: greille céanna" má cheapann tú gurb í an ghreille chéanna í nó "CINNEADH: greillí éagsúla" má cheapann tú gur greillí éagsúla iad. Ná cuir aon rud eile leis.} \\
            }
        }
    }
    & \\ \\

    \theutterance \stepcounter{utterance}  
    & & & \multicolumn{4}{p{0.6\linewidth}}{
        \cellcolor[rgb]{0.9,0.9,0.9}{
            \makecell[{{p{\linewidth}}}]{
                \texttt{\tiny{[GM$\langle$P2]}}
                \texttt{CINNEADH: greillí éagsúla} \\
            }
        }
    }
     \\ \\

    \theutterance \stepcounter{utterance}  
    & & & \multicolumn{2}{p{0.3\linewidth}}{
        \cellcolor[rgb]{0.9,0.9,0.9}{
            \makecell[{{p{\linewidth}}}]{
                \texttt{\tiny{[GM$|$GM]}}
                \texttt{continue} \\
            }
        }
    }
    & & \\ \\

    \theutterance \stepcounter{utterance}  
    & & & \multicolumn{2}{p{0.3\linewidth}}{
        \cellcolor[rgb]{0.9,0.9,0.9}{
            \makecell[{{p{\linewidth}}}]{
                \texttt{\tiny{[GM$|$GM]}}
                \texttt{success} \\
            }
        }
    }
    & & \\ \\

    \theutterance \stepcounter{utterance}  
    & & \multicolumn{4}{p{0.6\linewidth}}{
        \cellcolor[rgb]{0.9,0.9,0.9}{
            \makecell[{{p{\linewidth}}}]{
                \texttt{\tiny{[P1$\langle$GM]}}
                \texttt{FREAGRA: Ní hea. Sa tríú fhadlíne, tá 'X' sa chéad, sa tríú agus sa chúigiú ceallóg amháin. Tá an dara agus an ceathrú ceallóg folamh (▢).} \\
\texttt{Anois déan cinneadh. Cad a cheapann tú: an bhfuil do ghreille agus an ghreille eile a thuairiscítear mar an gcéanna? Scríobh "CINNEADH: greille céanna" má cheapann tú gurb í an ghreille chéanna í nó "CINNEADH: greillí éagsúla" má cheapann tú gur greillí éagsúla iad. Ná cuir aon rud eile leis.} \\
            }
        }
    }
    & \\ \\

    \theutterance \stepcounter{utterance}  
    & \multicolumn{4}{p{0.6\linewidth}}{
        \cellcolor[rgb]{0.9,0.9,0.9}{
            \makecell[{{p{\linewidth}}}]{
                \texttt{\tiny{[P1$\rangle$GM]}}
                \texttt{CINNEADH: greillí éagsúla} \\
            }
        }
    }
    & & \\ \\

    \theutterance \stepcounter{utterance}  
    & & & \multicolumn{2}{p{0.3\linewidth}}{
        \cellcolor[rgb]{0.9,0.9,0.9}{
            \makecell[{{p{\linewidth}}}]{
                \texttt{\tiny{[GM$|$GM]}}
                \texttt{continue} \\
            }
        }
    }
    & & \\ \\

    \theutterance \stepcounter{utterance}  
    & & & \multicolumn{2}{p{0.3\linewidth}}{
        \cellcolor[rgb]{0.9,0.9,0.9}{
            \makecell[{{p{\linewidth}}}]{
                \texttt{\tiny{[GM$|$GM]}}
                \texttt{success} \\
            }
        }
    }
    & & \\ \\

\end{supertabular}
}

\end{document}
